% Generated draft: requires mathematical and source review.
\documentclass[11pt]{article}
\usepackage[margin=25mm]{geometry}
\usepackage{fontspec}
\setmainfont{Latin Modern Roman}
\usepackage{amsmath,amssymb,mathtools,unicode-math}
\setmathfont{Latin Modern Math}
\usepackage{enumitem,needspace,xurl,hyperref}
\setlength{\parindent}{0pt}
\setlength{\parskip}{5pt}
\setlength{\emergencystretch}{3em}
\newcommand{\N}{\mathbb{N}}
\newcommand{\Z}{\mathbb{Z}}
\newcommand{\Q}{\mathbb{Q}}
\newcommand{\R}{\mathbb{R}}
\newcommand{\C}{\mathbb{C}}
\newcommand{\F}{\mathbb{F}}
\newcommand{\A}{\mathbb{A}}
\newcommand{\PP}{\mathbb{P}}
\newcommand{\E}{\mathbb{E}}
\newcommand{\eps}{\varepsilon}
\newcommand{\dd}{\,\mathrm{d}}
\begin{document}
\section*{UnsolvedMath ID: 800571}
\textbf{Source problem code:} catalogue entry
\subsection{Definitions and mathematical statement}
The supplied excerpt does not define the partition function \(Z\) or specify the lattice-gas model. Consequently, the following is a conditional formulation of the question's quantifiers, rather than a complete definition of its mathematical object.

Fix the parameter \(\Delta\). Let \(\mathcal G_\Delta\) stand for the source's intended class of graphs of maximum degree \(\Delta\). For \(G\in\mathcal G_\Delta\), write \(V(G)\) for its vertex set and use \(Z_G((w_x)_{x\in V(G)})\) as a placeholder for the unspecified partition function, with complex vertex parameters \(w_x\). Its equal-parameter specialization is denoted by
\[
Z_G^{\mathrm{eq}}(w)=Z_G((w)_{x\in V(G)}).
\]
Here a complex domain means a connected open subset of \(\C\).

\textbf{Equal-parameter version.} Determine the largest domain \(D\subseteq\C\), containing \(0\), such that
\[
\forall G\in\mathcal G_\Delta\;\forall w\in D,
\qquad Z_G^{\mathrm{eq}}(w)\ne0.
\]
Largest means containing every other domain with this property; existence of such a largest domain is part of the intended question and is not asserted here.

\textbf{Unequal-parameter version.} Seek domains \(D\subseteq\C\), containing \(0\), such that
\[
\forall G\in\mathcal G_\Delta\;\forall (w_x)_{x\in V(G)}\in D^{V(G)},
\qquad Z_G((w_x)_{x\in V(G)})\ne0.
\]
The source explicitly cautions that there need not be a unique maximal such domain. A domain is maximal here if no strictly larger domain containing it satisfies the same condition; this differs from being a largest domain.

A complete statement requires the source's definition of \(Z\), its graph conventions, the allowed values of \(\Delta\), and clarification of whether ``maximum degree \(\Delta\)'' means degree at most \(\Delta\) or maximum degree exactly \(\Delta\). No particular partition function or graph convention is supplied by this draft.
\subsection{Short English statement}
Sokal's Question 2 asks for the largest complex domain containing zero in which the partition function stays nonzero uniformly over the specified graphs. It distinguishes equal vertex parameters from independently varying vertex parameters; in the latter version, a unique maximal domain need not exist. The supplied excerpt leaves the partition function and graph conventions undefined, so the displayed formulation remains conditional.
\subsection{Sources}
\begin{itemize}
\item Local UnsolvedMath registry, permanent ID \(800571\), source problem code \texttt{catalogue entry}, supplied title: ``Lattice Gases and Antiferromagnetic Potts Models — Question 2.'' The registry attributes the problem to Sokal, proposed in 2000, and labels it \texttt{open}; this is a supplied status label, not an independently verified assessment.
\item Sokal, \emph{A Personal List of Unsolved Problems Concerning Lattice Gases and Antiferromagnetic Potts Models} (2000), Question 2. Supplied source URL: \url{https://arxiv.org/abs/cond-mat/0004231}. The supplied provenance records access on 2026-10-16 and extraction by \texttt{arxiv\_latex}; these are supplied provenance claims, not a claim of independent access or verification.
\item Supplied upstream revision: \texttt{d371ee917ac9647fbb06cd6fc529deec6f3dda82}.
\end{itemize}
\end{document}
